\chapter{Frontier Orbitals and Reactivity}\label{ch:b2:frontier-orbitals}

A bottle of cyclopentadiene left on the shelf slowly turns into its dimer:
two molecules of the same \omterm{def:b2:frontier-orbitals:diels-alder}{diene} join into a bridged bicyclic compound, and
only one of the possible stereoisomers forms. Nothing in the curly arrows of
the Year 1 volume predicts which. Two orbitals do: the highest occupied
orbital of one molecule and the lowest empty orbital of the other. This
chapter turns the two-level interaction of \cref{ch:b2:diatomic-mos} into a
tool for reactivity --- which atom a nucleophile attacks, from which side, and
why a \omterm{def:b2:frontier-orbitals:diels-alder}{diene} and an alkene combine into a ring in a single step.

\begin{recall}
\Cref{ch:b2:diatomic-mos}: two interacting orbitals, four-electron repulsion.
\Cref{ch:b2:huckel}: H\"uckel levels and coefficients of conjugated systems.
The Year 1 volume: nucleophiles and electrophiles, curly arrows, kinetic
control, stereospecific, stereoselective and regioselective reactions,
electrophilic additions to alkenes, nucleophilic substitution.
\end{recall}

\section{Two orbitals, revisited}

When two molecules approach, their orbitals overlap and interact. Most of the
interaction is weak, the orbitals being far apart in energy; a perturbative form of
the two-level result then suffices.

\begin{theorem}[Weak interaction of two levels]\label{thm:b2:frontier-orbitals:perturbation}
Two orbitals of energies $\varepsilon_1 < \varepsilon_2$, coupled by a \omterm{def:b2:diatomic-mos:integrals}{resonance integral}
$\beta$ with $|\beta| \ll \Delta\varepsilon = \varepsilon_2 - \varepsilon_1$ (overlap neglected), are
shifted to
\[
  \varepsilon_1 - \frac{\beta^2}{\Delta\varepsilon}, \qquad \varepsilon_2 + \frac{\beta^2}{\Delta\varepsilon}
\]
to first order in $\beta^2/\Delta\varepsilon^2$: the lower level falls and the upper rises by the same
amount, the smaller as the gap is larger.
\end{theorem}

\begin{proof}
By \cref{thm:b2:diatomic-mos:heteronuclear} the exact levels are $\bar\varepsilon \mp
\sqrt{\Delta\varepsilon^2/4 + \beta^2}$. With $u = 4\beta^2/\Delta\varepsilon^2 \ll 1$, $\sqrt{\Delta\varepsilon^2/4
+ \beta^2} = \frac{\Delta\varepsilon}{2}\sqrt{1 + u} \approx \frac{\Delta\varepsilon}{2}(1 + u/2) =
\frac{\Delta\varepsilon}{2} + \frac{\beta^2}{\Delta\varepsilon}$, and $\bar\varepsilon \mp \Delta\varepsilon/2 = \varepsilon_1,
\varepsilon_2$.
\end{proof}

\begin{proposition}[Two electrons attract, four repel]\label{prop:b2:frontier-orbitals:four-electron}
If the two interacting orbitals hold two electrons in all, the system is stabilised
by about $2\beta^2/\Delta\varepsilon$. If they hold four, the system is destabilised.
\end{proposition}

\begin{proof}
Two electrons go into the lowered level: gain $2\beta^2/\Delta\varepsilon$. With four, the upper level
is filled too; to the order of \cref{thm:b2:frontier-orbitals:perturbation} the two
shifts cancel, and with the overlap kept the upper level rises more than the lower
falls (\cref{prop:b2:diatomic-mos:asymmetry}): a net repulsion.
\end{proof}

\begin{omfigure}
\begin{tikzpicture}[font=\small]
  \begin{scope}
    \pic (a) at (0,0) {omlevel={ud}{0.7}};
    \pic (b) at (3,1.4) {omlevel={empty}{0.7}};
    \pic (lo) at (1.5,-0.5) {omlevel={ud}{0.7}};
    \pic (hi) at (1.5,1.9) {omlevel={empty}{0.7}};
    \draw[omcorr, densely dashed] (a-r) -- (lo-l) (a-r) -- (hi-l) (b-l) -- (lo-r) (b-l) -- (hi-r);
    \node[below=6pt] at (1.5,-0.6) {two electrons: stabilised};
    \node[left] at (a-l) {filled};
    \node[right] at (b-r) {empty};
  \end{scope}
  \begin{scope}[xshift=7cm]
    \pic (a) at (0,0) {omlevel={ud}{0.7}};
    \pic (b) at (3,0.6) {omlevel={ud}{0.7}};
    \pic (lo) at (1.5,-0.5) {omlevel={ud}{0.7}};
    \pic (hi) at (1.5,1.25) {omlevel={ud}{0.7}};
    \draw[omcorr, densely dashed] (a-r) -- (lo-l) (a-r) -- (hi-l) (b-l) -- (lo-r) (b-l) -- (hi-r);
    \node[below=6pt] at (1.5,-0.6) {four electrons: repelled};
    \node[left] at (a-l) {filled};
    \node[right] at (b-r) {filled};
  \end{scope}
\end{tikzpicture}

{\small Interaction of two orbitals of neighbouring molecules. A filled orbital
facing an empty one gives a net stabilisation; two filled orbitals repel, the
upper level rising more than the lower one falls. This repulsion is the origin
of steric hindrance between molecules.}
\end{omfigure}

\section{Frontier orbitals}

\begin{definition}[Frontier orbitals]\label{def:b2:frontier-orbitals:frontier}
The \emph{HOMO}\index{HOMO} (highest occupied \omterm{def:b2:diatomic-mos:lcao}{molecular orbital}) and the
\emph{LUMO}\index{LUMO} (lowest unoccupied \omterm{def:b2:diatomic-mos:lcao}{molecular orbital}) of a molecule are its
\emph{frontier orbitals}\index{frontier orbitals}.
\end{definition}

\begin{proposition}[The dominant interaction]\label{prop:b2:frontier-orbitals:dominant}
Between two closed-shell molecules, the stabilising interactions that matter most
are between the \omterm{def:b2:frontier-orbitals:frontier}{HOMO} of one and the \omterm{def:b2:frontier-orbitals:frontier}{LUMO} of the other; of the two such pairs, the one
with the smaller gap dominates. A nucleophile reacts through its \omterm{def:b2:frontier-orbitals:frontier}{HOMO}, an
electrophile through its \omterm{def:b2:frontier-orbitals:frontier}{LUMO}.
\end{proposition}

\begin{proof}[Argument]
Filled--filled pairs repel (\cref{prop:b2:frontier-orbitals:four-electron}) and empty--empty
pairs hold no electrons: only filled--empty pairs stabilise, each by $2\beta^2/\Delta\varepsilon$.
The \omterm{def:b2:frontier-orbitals:frontier}{HOMO} is the highest filled level and the \omterm{def:b2:frontier-orbitals:frontier}{LUMO} the lowest empty one, so a
HOMO--LUMO pair has the smallest gap among them; by the theorem, the smallest gap
gives the largest term.
\end{proof}

\begin{definition}[Charge control and orbital control]\label{def:b2:frontier-orbitals:control}
A reaction is under \emph{charge control}\index{charge control} when the attraction
between the charges of the two partners decides where they react, and under
\emph{orbital control}\index{orbital control} when the overlap of the \omterm{def:b2:frontier-orbitals:frontier}{frontier orbitals}
decides it, the attack taking place where their coefficients are largest.
\end{definition}

\begin{definition}[Hard and soft species]\label{def:b2:frontier-orbitals:hard-soft}
A \emph{hard nucleophile}\index{hard nucleophile} or \emph{hard
electrophile}\index{hard electrophile} is small, carries a concentrated charge and has
a \omterm{def:b2:frontier-orbitals:frontier}{HOMO} (respectively a \omterm{def:b2:frontier-orbitals:frontier}{LUMO}) far from that of its partners: it reacts under \omterm{def:b2:frontier-orbitals:control}{charge
control}. A \emph{soft nucleophile}\index{soft nucleophile} or \emph{soft
electrophile}\index{soft electrophile} is large and polarisable, with a high \omterm{def:b2:frontier-orbitals:frontier}{HOMO}
(respectively a low \omterm{def:b2:frontier-orbitals:frontier}{LUMO}): it reacts under \omterm{def:b2:frontier-orbitals:control}{orbital control}. Hard reacts preferably
with hard, soft with soft.
\end{definition}

The hydroxide ion and the fluoride ion are \omterm{def:b2:frontier-orbitals:hard-soft}{hard nucleophiles}, iodide and thiolates
soft; a proton and a carbocation are \omterm{def:b2:frontier-orbitals:hard-soft}{hard electrophiles}, the carbon of a
halogenoalkane soft.

\begin{definition}[Ambident nucleophile]\label{def:b2:frontier-orbitals:ambident}
An \emph{ambident nucleophile}\index{ambident nucleophile} has two nucleophilic atoms
connected by conjugation, which can each attack an electrophile.
\end{definition}

The cyanide ion is attacked at nitrogen, where the negative charge is larger, by \omterm{def:b2:frontier-orbitals:hard-soft}{hard
electrophiles}, and at carbon, where its \omterm{def:b2:frontier-orbitals:frontier}{HOMO} has its larger coefficient, by soft ones:
a halogenoalkane gives a nitrile \ce{R-C#N}. Enolate ions, met in
\cref{ch:b2:enolates-aldol}, react at carbon or at oxygen in the same way.

\begin{method}[A frontier-orbital analysis]\label{met:b2:frontier-orbitals:fmo}
\begin{enumerate}
  \item Identify the nucleophile (high \omterm{def:b2:frontier-orbitals:frontier}{HOMO}) and the electrophile (low \omterm{def:b2:frontier-orbitals:frontier}{LUMO}).
  \item Draw the two \omterm{def:b2:frontier-orbitals:frontier}{frontier orbitals} with their coefficients and signs.
  \item The reaction occurs where the overlap is largest: atoms with the largest
        coefficients, lobes of the same sign facing each other.
  \item The direction of approach follows the shape of the \omterm{def:b2:frontier-orbitals:frontier}{LUMO} (perpendicular to a
        \ce{C=O} plane, at the back of a \ce{C-X} bond).
  \item A smaller HOMO--LUMO gap means a faster reaction (other things being equal).
\end{enumerate}
\end{method}

The direction of attack follows from the \omterm{def:b2:frontier-orbitals:frontier}{LUMO}. The \omterm{def:b2:frontier-orbitals:frontier}{LUMO} of a carbonyl group is its
$\pi^*$ orbital, larger on carbon and perpendicular to the plane: nucleophiles approach
the carbon from above or below the plane, not along the \ce{C=O} axis. The \omterm{def:b2:frontier-orbitals:frontier}{LUMO} of a
halogenoalkane is the $\sigma^*$ orbital of the \ce{C-X} bond, whose large lobe on carbon
points away from X: the nucleophile attacks from the back, and the carbon is inverted
--- the stereochemistry of the bimolecular substitution of the Year 1 volume.

\section{The Diels--Alder reaction}

\begin{definition}[Concerted reaction]\label{def:b2:frontier-orbitals:concerted}
A \emph{concerted reaction}\index{concerted reaction} makes and breaks all its bonds in
a single step, through one transition state, without an intermediate.
\end{definition}

\begin{definition}[Diels--Alder reaction]\label{def:b2:frontier-orbitals:diels-alder}
The \emph{Diels--Alder reaction}\index{Diels--Alder reaction} is the concerted addition
of a conjugated \emph{diene}\index{diene}, in its s-cis conformation, to an alkene or
alkyne, the \emph{dienophile}\index{dienophile}, forming a six-membered ring with two
new $\sigma$ bonds and one new $\pi$ bond.
\end{definition}

\begin{omfigure}
\begin{tikzpicture}[font=\small, >={Stealth}, scale=1.25]
  % diene s-cis: C1 (0,1.6) C2 (0.6,2.4) C3 (1.6,2.4) C4 (2.2,1.6); dienophile below
  \draw[thick] (0,1.6) -- (0.6,2.4) -- (1.6,2.4) -- (2.2,1.6);
  \draw[thick] (0.2,1.62) -- (0.65,2.22);
  \draw[thick] (2.0,1.62) -- (1.55,2.22);
  \draw[thick] (0.2,0) -- (2.0,0); \draw[thick] (0.35,0.14) -- (1.85,0.14);
  \node[font=\scriptsize, anchor=east] at (-0.05,1.6) {1};
  \node[font=\scriptsize, anchor=south east] at (0.5,2.42) {2};
  \node[font=\scriptsize, anchor=south] at (1.6,2.42) {3};
  \node[font=\scriptsize, anchor=west] at (2.25,1.6) {4};
  % arrows: dienophile pi -> C1..C(dienophile); C1=C2 pi -> C2-C3; C3=C4 pi -> C4-C
  \draw[->, thick, omProp] (1.1,0.22) .. controls (0.9,0.9) and (0.3,0.9) .. (0.12,1.42);
  \draw[->, thick, omProp] (0.45,1.95) .. controls (0.6,2.75) and (1.0,2.8) .. (1.1,2.48);
  \draw[->, thick, omProp] (1.78,1.92) .. controls (2.5,1.6) and (2.4,0.6) .. (2.05,0.2);
  \node at (3.3,1.2) {$\longrightarrow$};
  \begin{scope}[xshift=4.2cm]
    \draw[thick] (0,1.6) -- (0.6,2.4) -- (1.6,2.4) -- (2.2,1.6) -- (2.0,0) -- (0.2,0) -- cycle;
    \draw[thick] (0.72,2.26) -- (1.48,2.26);
  \end{scope}
\end{tikzpicture}

{\small The \omterm{def:b2:frontier-orbitals:diels-alder}{Diels--Alder reaction} of butadiene and ethene, written with curly arrows:
three pairs of electrons move at once, in a cyclic, concerted step; two $\sigma$ bonds
form at the ends of the \omterm{def:b2:frontier-orbitals:diels-alder}{diene}, and the $\pi$ bond moves to its centre.}
\end{omfigure}

The dominant interaction is between the \omterm{def:b2:frontier-orbitals:frontier}{HOMO} of the \omterm{def:b2:frontier-orbitals:diels-alder}{diene}, $\psi_2$, and the \omterm{def:b2:frontier-orbitals:frontier}{LUMO} of the
\omterm{def:b2:frontier-orbitals:diels-alder}{dienophile}, $\pi^*$. Their lobes at the ends of the \omterm{def:b2:frontier-orbitals:diels-alder}{diene} and on the two carbons of the
\omterm{def:b2:frontier-orbitals:diels-alder}{dienophile} have matching signs when the \omterm{def:b2:frontier-orbitals:diels-alder}{dienophile} approaches the face of the \omterm{def:b2:frontier-orbitals:diels-alder}{diene}:
both new bonds can form at once, on the same face of each partner.

\begin{omfigure}
\begin{tikzpicture}[font=\small]
  \foreach \x/\c in {0/0.602, 1/0.372, 2/-0.372, 3/-0.602} {
    \pic at (\x,2) {omp={1.5*\c}{90}};
  }
  \draw[black!45] (-0.2,2) -- (3.2,2);
  \node[anchor=west] at (3.6,2) {HOMO of the diene ($\psi_2$)};
  \pic at (0,0) {omp={-0.85}{90}}; \pic at (3,0) {omp={0.85}{90}};
  \draw[black!45] (-0.2,0) -- (3.2,0);
  \node[anchor=west] at (3.6,0) {LUMO of the dienophile ($\pi^*$)};
  \draw[omProp, thick, densely dashed] (0,1.45) -- (0,0.6);
  \draw[omProp, thick, densely dashed] (3,1.45) -- (3,0.6);
  \node[font=\scriptsize, omProp, anchor=east] at (-0.15,1.0) {in phase};
  \node[font=\scriptsize, omProp, anchor=west] at (3.15,1.0) {in phase};
\end{tikzpicture}

{\small \omterm{def:b2:frontier-orbitals:frontier}{Frontier orbitals} of butadiene and ethene, lobes in proportion to their
H\"uckel coefficients. Facing each other, the end lobes of the \omterm{def:b2:frontier-orbitals:diels-alder}{diene}'s \omterm{def:b2:frontier-orbitals:frontier}{HOMO} and
the lobes of the \omterm{def:b2:frontier-orbitals:diels-alder}{dienophile}'s \omterm{def:b2:frontier-orbitals:frontier}{LUMO} have the same signs at both ends: the two
$\sigma$ bonds can form together, on the same face of each partner.}
\end{omfigure}

\begin{proposition}[Stereospecificity]\label{prop:b2:frontier-orbitals:da-stereospecific}
The \omterm{def:b2:frontier-orbitals:diels-alder}{Diels--Alder reaction} is stereospecific: substituents cis on the \omterm{def:b2:frontier-orbitals:diels-alder}{dienophile} are
cis in the product, trans substituents remain trans.
\end{proposition}

\begin{proof}
The reaction is concerted and both new bonds form on the same face of the
\omterm{def:b2:frontier-orbitals:diels-alder}{dienophile}: its two carbons never rotate with respect to each other, and the relative
positions of their substituents are carried into the ring.
\end{proof}

\begin{proposition}[Regioselectivity]\label{prop:b2:frontier-orbitals:da-regio}
With unsymmetrical partners, the major product joins the atom of the \omterm{def:b2:frontier-orbitals:diels-alder}{diene} with the
largest coefficient in its \omterm{def:b2:frontier-orbitals:frontier}{HOMO} to the atom of the \omterm{def:b2:frontier-orbitals:diels-alder}{dienophile} with the largest
coefficient in its \omterm{def:b2:frontier-orbitals:frontier}{LUMO}.
\end{proposition}

\begin{proof}[Argument]
In the transition state the two new bonds are not yet equally formed; the
stabilisation $2\beta^2/\Delta\varepsilon$ is larger when the larger coefficients overlap ($\beta$
grows with the product of the coefficients of the overlapping atoms), so that
orientation is reached at lower energy.
\end{proof}

\begin{omfigure}
\begin{tikzpicture}
\begin{axis}[ybar, width=0.47\linewidth, height=4.6cm, ymin=-0.8, ymax=0.8,
  xtick={0,1,2,3,4}, xticklabels={O,C1,C2,C3,C4}, axis x line=bottom, axis y line=left,
  extra y ticks={0}, extra y tick style={grid=major, grid style={black!50}}, enlarge x limits=0.15,
  tick label style={font=\scriptsize}, ylabel={coefficient}, label style={font=\small},
  title={\small HOMO of 1-methoxybutadiene}, bar width=8pt]
\addplot[fill=omDef!50, draw=omDef] table[x=atom, y=c] {figdata/out/bachelor-2-huckel-c-diene.dat};
\end{axis}
\end{tikzpicture}\hfill
\begin{tikzpicture}
\begin{axis}[ybar, width=0.47\linewidth, height=4.6cm, ymin=-0.8, ymax=0.8,
  xtick={1,2,3,4}, xticklabels={C$\beta$,C$\alpha$,C(O),O}, axis x line=bottom, axis y line=left,
  extra y ticks={0}, extra y tick style={grid=major, grid style={black!50}}, enlarge x limits=0.15,
  tick label style={font=\scriptsize}, ylabel={coefficient}, label style={font=\small},
  title={\small LUMO of propenal}, bar width=8pt]
\addplot[fill=omThm!45, draw=omThm] table[x=atom, y=c] {figdata/out/bachelor-2-huckel-c-dienophile.dat};
\end{axis}
\end{tikzpicture}

{\small H\"uckel model with heteroatom parameters ($\alpha_{\ce{O}} = \alpha + 2\beta$,
$\beta_{\ce{CO}} = 0.8\beta$ for the ether oxygen; $\alpha_{\ce{O}} = \alpha + \beta$, $\beta_{\ce{CO}} =
\beta$ for the carbonyl oxygen; a model). The \omterm{def:b2:frontier-orbitals:diels-alder}{diene}'s \omterm{def:b2:frontier-orbitals:frontier}{HOMO} is largest on C4, the end
far from the methoxy group (0.60 against 0.51); the \omterm{def:b2:frontier-orbitals:diels-alder}{dienophile}'s \omterm{def:b2:frontier-orbitals:frontier}{LUMO} is largest on
the $\beta$ carbon (0.66). They bond together: the methoxy and aldehyde groups end
up on neighbouring carbons of the ring.}
\end{omfigure}

The donor group raises the \omterm{def:b2:frontier-orbitals:diels-alder}{diene}'s \omterm{def:b2:frontier-orbitals:frontier}{HOMO} ($m = 0.48$ instead of 0.62) and the acceptor
lowers the \omterm{def:b2:frontier-orbitals:diels-alder}{dienophile}'s \omterm{def:b2:frontier-orbitals:frontier}{LUMO} ($m = -0.35$ instead of $-1$): the gap narrows, and such
pairs react much faster than butadiene with ethene. A \omterm{def:b2:frontier-orbitals:diels-alder}{diene} rich in electrons and a
\omterm{def:b2:frontier-orbitals:diels-alder}{dienophile} poor in electrons is the typical combination.

\begin{definition}[Endo and exo adducts]\label{def:b2:frontier-orbitals:endo}
When a cyclic \omterm{def:b2:frontier-orbitals:diels-alder}{diene} adds to a \omterm{def:b2:frontier-orbitals:diels-alder}{dienophile} carrying an unsaturated substituent, the
\emph{endo adduct}\index{endo adduct} has that substituent pointing towards the newly
formed double bond (under the \omterm{def:b2:frontier-orbitals:diels-alder}{diene} in the transition state), the \emph{exo
adduct}\index{exo adduct} away from it. The \emph{endo rule}\index{endo rule} states
that the endo adduct is formed faster.
\end{definition}

\begin{proposition}[The endo rule]\label{prop:b2:frontier-orbitals:endo}
Under kinetic control, the \omterm{def:b2:frontier-orbitals:endo}{endo adduct} is the major product, although the \omterm{def:b2:frontier-orbitals:endo}{exo adduct}
is often the more stable.
\end{proposition}

\admitted

The usual explanation is a secondary overlap, in the endo approach, between the
lobes of the substituent's $\pi$ system and those of the central atoms of the \omterm{def:b2:frontier-orbitals:diels-alder}{diene}'s
\omterm{def:b2:frontier-orbitals:frontier}{HOMO}; it lowers the endo transition state without forming a bond. The rule is
empirical and has exceptions; the Year 3 volume treats these reactions with the
symmetry of their orbitals.

\begin{omfigure}
\begin{tikzpicture}[font=\small]
  \foreach \xs/\mode in {0/endo, 5/exo} {
    \begin{scope}[xshift=\xs cm]
      % cyclopentadiene seen from the side: C1, C2=C3, C4, CH2 bridge on top
      \draw[thick] (0,2) -- (0.6,2.5) -- (1.8,2.5) -- (2.4,2);
      \draw[thick] (0.75,2.38) -- (1.65,2.38);
      \draw[thick] (0,2) -- (1.2,3.2) -- (2.4,2);
      \node[font=\scriptsize, anchor=south] at (1.2,3.2) {\ce{CH2}};
      % dienophile
      \draw[thick] (0.2,0) -- (2.2,0); \draw[thick] (0.35,0.1) -- (2.05,0.1);
      \draw[omProp, densely dashed, thick] (0,1.9) -- (0.2,0.15);
      \draw[omProp, densely dashed, thick] (2.4,1.9) -- (2.2,0.15);
      \node[anchor=north] at (1.2,-1.3) {\mode};
    \end{scope}
  }
  % endo: carbonyls tucked under the diene
  \draw[thick] (0.2,0) -- (0.75,0.95) (2.2,0) -- (1.65,0.95);
  \node[font=\scriptsize, fill=white, inner sep=1pt] at (0.75,1.1) {C=O};
  \node[font=\scriptsize, fill=white, inner sep=1pt] at (1.65,1.1) {C=O};
  \draw[omThm, densely dotted, thick] (0.75,1.25) -- (0.65,2.35);
  \draw[omThm, densely dotted, thick] (1.65,1.25) -- (1.75,2.35);
  % exo: carbonyls pointing away
  \draw[thick] (5.2,0) -- (4.95,-0.85) (7.2,0) -- (7.45,-0.85);
  \node[font=\scriptsize] at (4.95,-1.0) {C=O};
  \node[font=\scriptsize] at (7.45,-1.0) {C=O};
\end{tikzpicture}

{\small Endo and exo approaches of maleic anhydride to cyclopentadiene, schematic.
The forming bonds are dashed. In the endo approach the carbonyl groups lie under the
\omterm{def:b2:frontier-orbitals:diels-alder}{diene}, and their $\pi$ orbitals overlap with its central carbons (dotted, secondary
overlap); in the exo approach they point away.}
\end{omfigure}

\begin{method}[Drawing a Diels--Alder product]\label{met:b2:frontier-orbitals:diels-alder}
\begin{enumerate}
  \item Put the \omterm{def:b2:frontier-orbitals:diels-alder}{diene} in its s-cis conformation, the \omterm{def:b2:frontier-orbitals:diels-alder}{dienophile} facing its ends.
  \item Draw the two new $\sigma$ bonds from the \omterm{def:b2:frontier-orbitals:diels-alder}{diene}'s ends to the \omterm{def:b2:frontier-orbitals:diels-alder}{dienophile}'s carbons,
        and the new $\pi$ bond between C2 and C3 of the \omterm{def:b2:frontier-orbitals:diels-alder}{diene}.
  \item Keep the \omterm{def:b2:frontier-orbitals:diels-alder}{dienophile}'s substituents on the same side (cis stays cis).
  \item For unsymmetrical partners, join the largest coefficients (``ortho'' or ``para''
        products for 1- or 2-substituted \omterm{def:b2:frontier-orbitals:diels-alder}{dienes}).
  \item With a cyclic \omterm{def:b2:frontier-orbitals:diels-alder}{diene}, put the unsaturated substituent endo.
\end{enumerate}
\end{method}

\begin{inthelab}[Cracking dicyclopentadiene]
Cyclopentadiene is sold as its dimer and regenerated just before use: the dimer is
heated to its boiling point (about \qty{170}{\degreeCelsius}) in a flask fitted with a
fractionating column, where it splits back into the monomer (a retro-Diels--Alder
reaction), which distils at \qty{41}{\degreeCelsius} and is collected in a receiver cooled
in ice. The monomer dimerises again within hours at room temperature and is used at
once. Both compounds are flammable and harmful; the work is done in a fume hood.
\end{inthelab}

\begin{safety}{\ghs[5mm]{GHS02}\,\ghs[5mm]{GHS06}\,\ghs[5mm]{GHS08}}
Cyclopentadiene and its dimer are highly flammable and toxic by inhalation; maleic
anhydride is corrosive and a respiratory sensitiser. They are handled in a fume hood,
with gloves and eye protection, away from flames.
% ledger: ghs:cyclopentadiene, ghs:dicyclopentadiene, ghs:maleic-anhydride, bp:cyclopentadiene, bp:dicyclopentadiene
\end{safety}

\section{Exercises}

\begin{exercise}[$\star$]\label{exo:b2:frontier-orbitals:1}
Give the \omterm{def:b2:frontier-orbitals:frontier}{HOMO} and the \omterm{def:b2:frontier-orbitals:frontier}{LUMO} (energies in H\"uckel units) of ethene, of the allyl anion
and of butadiene.
\end{exercise}

\begin{exercise}[$\star$]\label{exo:b2:frontier-orbitals:2}
Classify as hard or soft: \ce{OH-}, \ce{I-}, \ce{H+}, a thiolate \ce{RS-}, the carbon of
iodomethane, \ce{F-}.
\end{exercise}

\begin{exercise}[$\star$]\label{exo:b2:frontier-orbitals:3}
Which of these \omterm{def:b2:frontier-orbitals:diels-alder}{dienes} can take part in a \omterm{def:b2:frontier-orbitals:diels-alder}{Diels--Alder reaction}: buta-1,3-diene,
cyclopentadiene, penta-1,4-diene, (2\textit{Z},4\textit{Z})-hexa-2,4-diene in its s-cis
form?
\end{exercise}

\begin{exercise}[$\star$]\label{exo:b2:frontier-orbitals:4}
Draw the product of butadiene with propenenitrile \ce{CH2=CH-C#N}.
\end{exercise}

\begin{exercise}[$\star\star$]\label{exo:b2:frontier-orbitals:5}
Two orbitals at $-10$ and \qty{-2}{eV} interact with $\beta = \qty{-1.0}{eV}$ (exercise data).
Compute the perturbative stabilisation of two electrons, then the exact one.
\end{exercise}

\begin{exercise}[$\star\star$]\label{exo:b2:frontier-orbitals:6}
The cyanide ion reacts with iodomethane to give ethanenitrile \ce{CH3-C#N}, but with a
carbocation in a strongly ionising solvent it gives some isonitrile \ce{R-N#C}. Explain.
\end{exercise}

\begin{exercise}[$\star\star$]\label{exo:b2:frontier-orbitals:7}
Using the coefficients of the chapter, predict the major product of 1-methoxybutadiene
with propenal.
\end{exercise}

\begin{exercise}[$\star\star$]\label{exo:b2:frontier-orbitals:8}
Draw the \omterm{def:b2:frontier-orbitals:endo}{endo adduct} of cyclopentadiene with methyl propenoate.
\end{exercise}

\begin{exercise}[$\star\star$]\label{exo:b2:frontier-orbitals:9}
Why does maleic anhydride react with cyclopentadiene at room temperature, while ethene
needs heating under pressure?
\end{exercise}

\begin{exercise}[$\star\star\star$]\label{exo:b2:frontier-orbitals:10}
Cyclopentadiene reacts with dimethyl maleate (cis diester) and with dimethyl fumarate
(trans diester). Draw both products and say how many stereoisomers each gives.
\end{exercise}

\begin{exercise}[$\star\star\star$]\label{exo:b2:frontier-orbitals:11}
The \omterm{def:b2:frontier-orbitals:diels-alder}{Diels--Alder reaction} has $\Delta_r H^\circ < 0$ and $\Delta_r S^\circ < 0$. Explain the signs
and why the reverse reaction becomes favourable at high temperature.
\end{exercise}

\begin{exercise}[$\star\star\star$]\label{exo:b2:frontier-orbitals:12}
In hydrazine \ce{H2N-NH2}, each nitrogen carries a lone pair. Using the four-electron
interaction, explain why the molecule avoids the conformation in which the two lone
pairs are parallel.
\end{exercise}

\section{Problem: Cracking Cyclopentadiene}

\begin{problem}[{Weekend problem --- the dimerisation of cyclopentadiene as a
Diels--Alder reaction, the thermodynamics of cracking, the frontier orbitals of
cyclopentadiene and maleic anhydride, and the stereochemistry of their adduct}]
\label{pb:b2:frontier-orbitals:1}
Data: boiling points, cyclopentadiene \qty{41}{\degreeCelsius}, dicyclopentadiene about
\qty{170}{\degreeCelsius}. H\"uckel model energies (a model): cyclopentadiene, treated as
a butadiene unit, \omterm{def:b2:frontier-orbitals:frontier}{HOMO} $\alpha + 0.62\beta$ and \omterm{def:b2:frontier-orbitals:frontier}{LUMO} $\alpha - 0.62\beta$; maleic anhydride, \omterm{def:b2:frontier-orbitals:frontier}{LUMO}
about $\alpha - 0.35\beta$.
% ledger: bp:cyclopentadiene, bp:dicyclopentadiene

\medskip
\noindent\textbf{Part I --- The dimer.}
\begin{enumerate}
  \item Draw cyclopentadiene and show that its \omterm{def:b2:frontier-orbitals:diels-alder}{diene} unit is held s-cis.
  \item In the dimerisation, one molecule is the \omterm{def:b2:frontier-orbitals:diels-alder}{diene}. What is the other?
  \item Draw the dimer and number the new bonds.
  \item Which adduct, endo or exo, forms at room temperature, and why?
  \item Is the dimerisation stereospecific? Explain with the mechanism.
  \item Why does the dimerisation not need a catalyst?
\end{enumerate}

\medskip
\noindent\textbf{Part II --- Thermodynamics of cracking.}
\begin{enumerate}[resume]
  \item Give the sign of $\Delta_r H^\circ$ of the dimerisation, from the bonds formed and
        broken.
  \item Give the sign of $\Delta_r S^\circ$.
  \item Write $\Delta_r G^\circ(T)$ and show that it changes sign at a temperature $T_i$.
  \item Why is the dimer stable at room temperature, and the monomer favoured when hot?
  \item In the cracking set-up, why does distilling the monomer as it forms drive the
        reaction?
  \item Why must the monomer be kept cold and used quickly?
\end{enumerate}

\medskip
\noindent\textbf{Part III --- \omterm{def:b2:frontier-orbitals:frontier}{Frontier orbitals}.}
\begin{enumerate}[resume]
  \item For cyclopentadiene reacting with itself, compute the HOMO--LUMO gap in units of
        $|\beta|$.
  \item For cyclopentadiene with maleic anhydride, compute the gap between the \omterm{def:b2:frontier-orbitals:diels-alder}{diene}'s
        \omterm{def:b2:frontier-orbitals:frontier}{HOMO} and the anhydride's \omterm{def:b2:frontier-orbitals:frontier}{LUMO}.
  \item Which pair reacts faster? Explain.
  \item Why do the two carbonyl groups of maleic anhydride lower its \omterm{def:b2:frontier-orbitals:frontier}{LUMO}?
  \item Which orbital of maleic anhydride gives the secondary overlap of the endo
        approach?
\end{enumerate}

\medskip
\noindent\textbf{Part IV --- The adduct with maleic anhydride.}
\begin{enumerate}[resume]
  \item Draw the \omterm{def:b2:frontier-orbitals:endo}{endo adduct}.
  \item Mark its stereocentres.
  \item Is the adduct chiral? (Look for a mirror plane.)
  \item Why do the two hydrogens of the former \omterm{def:b2:frontier-orbitals:diels-alder}{dienophile} end up on the same side of the
        ring?
  \item What would the \omterm{def:b2:frontier-orbitals:endo}{exo adduct} look like?
  \item Can the endo and \omterm{def:b2:frontier-orbitals:endo}{exo adducts} interconvert at room temperature? Why?
  \item How would the anhydride react with water, and what product would form?
  \item State the number of stereocentres of the \omterm{def:b2:frontier-orbitals:endo}{endo adduct} of cyclopentadiene and
        maleic anhydride.
\end{enumerate}
\end{problem}
